\chapter{Partial-PSUM 反馈与系统集成}
\label{chap:psum}

多 K pass 卷积的核心难点在于 partial-PSUM feedback。对于 $K_{total}>K_{tile}$ 的层，第一轮 K pass 只能产生部分和，后续 pass 必须读取前一轮 partial sum 并继续累加，最后一轮才进入重量化和输出路径。若每轮 pass 都严格串行执行“计算完成、全部 drain、下一轮 feed、下一轮 compute”，后端大层会出现明显阵列空泡。本章介绍本文如何对 partial-PSUM 路径和 pass 边界进行流水化与重叠。

\section{Partial-PSUM 反馈语义}

对一个输出 pixel 和一个 $Cout_{tile}$ 通道块，硬件内部形成一个 PSUM packet：
\begin{equation}
P[p,c] = \sum_{k=k_0}^{k_0+K_{tile}-1} X[p,k]\cdot W[k,c].
\end{equation}
当当前 pass 不是 final pass 时，该 packet 被写入片上 partial-PSUM buffer；下一 pass 读取同一地址的 partial sum 作为初值继续累加。当当前 pass 是 final pass 时，该 packet 被送入 requant/activation/pool/OFM path。

该语义要求三个条件同时成立：第一，partial-PSUM 的地址和 Cout block 必须与下一 pass 对齐；第二，写入和读取不能发生读写冲突；第三，final pass 的输出顺序必须与软件重排和 golden 比对一致。本文使用 ping-pong bank 和 per-pass context 维护这些条件。

\section{流水化 PSUM Drain}

早期 drain writer 在下游 ready 出现回压时容易产生地址推进和 valid 保持之间的时序风险，同时吞吐受内部状态机限制。本文将 drain writer 改为 AXI-Stream 风格握手：packet valid、addr 和 data 在真实 $valid \& ready$ 握手前保持稳定，握手成功后才推进地址。随后加入一拍 read-return tracker 和 one-entry hold register，使下游 ready 时可以接近每周期输出一个 PSUM packet。

流水化 drain 的意义不仅是提高 drain 吞吐，也为后续 overlap 提供可控事件：调度器可以基于 drain packet fire、drain done 和 FIFO empty wait 判断 partial-PSUM 是否已经产生足够可读数据。

\section{Early Drain 与 Continuous PSUM Collector}

传统 pass 调度在 compute done 后才启动 drain。Early drain 允许调度器在当前 pass 已经开始产生 PSUM 数据后提前启动 drain，只要 partial-PSUM FIFO 已经具备可读数据。该机制降低 compute 完成后等待 drain 启动的空泡，但仍保留 final pass 和非 final pass 的输出顺序约束。

Continuous PSUM collector 进一步改变非 final pass 的路径。collector 在 compute 开始前接收一个 pass context，并在阵列列 FIFO 产生完整 packet 后立即将其写入 partial-PSUM bank；final pass packet 则进入后处理流水线。这样，非 final pass 的 partial-PSUM 写回从“compute 后集中 drain”变为“compute 过程中持续收集”，减少 pass 边界处的反馈延迟。

图~\ref{fig:psum_prefetch_timeline} 对比了串行 pass 和重叠 pass 的时间线。

\begin{figure}[htb]
\centering
\resizebox{\textwidth}{!}{%
\begin{tikzpicture}[
    x=8mm,y=7mm,
    phase/.style={draw, minimum height=5mm, align=center},
    arr/.style={-Latex, thick}
]
\node at (0,2) {串行};
\node[phase, minimum width=14mm] at (2,2) {feed};
\node[phase, minimum width=18mm] at (4.1,2) {compute};
\node[phase, minimum width=18mm] at (6.4,2) {drain};
\node[phase, minimum width=16mm] at (8.5,2) {next feed};
\node[phase, minimum width=20mm] at (10.8,2) {next compute};

\node at (0,0) {重叠};
\node[phase, minimum width=14mm] at (2,0) {feed};
\node[phase, minimum width=28mm] at (4.6,0) {compute};
\node[phase, minimum width=24mm] at (5.3,-0.8) {continuous collect};
\node[phase, minimum width=24mm] at (7.6,0.8) {next prefetch};
\node[phase, minimum width=22mm] at (9.5,0) {next compute};
\draw[arr] (6.9,0.8) -- (8.5,0.1);
\end{tikzpicture}
}
\caption{Partial-PSUM 反馈与 pass-level prefetch 时间线}
\label{fig:psum_prefetch_timeline}
\end{figure}

\section{Pass-Level Prefetch}

K pass 之间除了 PSUM 反馈，还包含下一 pass 权重加载和 IFM replay。Pass-level prefetch 在当前 pass 已经满足安全条件后提前准备下一 pass 的 weight 和 raw-HWC replay。该机制遵循两个约束：一是不能覆盖当前 pass 正在使用的 PE 权重；二是 next-pass IFM replay 只能进入解耦 FIFO 或安全队列，不能破坏当前 pass 的输入顺序。

本文实现了两级 prefetch。第一阶段 pass prefetch 主要在当前 compute 完成后、drain 或 pass 边界窗口中准备下一 pass；第二阶段 during-compute prefetch 进一步允许在当前 pass compute 进行时启动下一 pass 权重和 raw-HWC replay 的 staging。后者能够降低 compute stage 中由于下一 pass 准备不足造成的 idle，但需要更严格地维护权重读取预算和 IFM FIFO 顺序。

\section{STREAM\_CFG 与调度开关}

所有优化均通过 AXI-Lite 中的 STREAM\_CFG 位控制。bit0 使能 batch mode，bit1 使能 raw-HWC mode，bit2 使能 early drain，bit3 使能 pass prefetch，bit4 使能 partial-PSUM overlap，bit5 使能 continuous PSUM collector，bit7 使能 during-compute prefetch。bit6 对应 column-level PSUM，是当前探索性设计，不作为最终推荐配置的必要条件。

这种设计保持了软件 ABI 的稳定性。默认关闭实验开关时，系统回到保守串行语义；逐步打开开关时，软件只需写入不同 STREAM\_CFG 和少量 tail/prefetch 参数，外部 AXI packet 格式保持不变。该策略便于在同一 runtime 中进行 A/B 测试和板级回归。

\section{性能计数器设计}

为了避免仅凭端到端延迟判断优化效果，本文在硬件中加入多层性能计数器。

\begin{table}[htb]
\centering
\caption{主要性能计数器类别}
\label{tab:perf_counters}
\begin{tabular}{ll}
\toprule
类别 & 作用 \\
\midrule
PERF & 统计 tile busy、外部等待、compute fire 等顶层周期 \\
STAGEPERF & 统计 bias、weight、feeder、compute、drain、OFM post 阶段 \\
SUBPERF & 细分 feeder push、FIFO stall、compute active、IFM stall、tail 等 \\
RAWSTAT & 统计 raw-HWC load、unpack、replay 和 replay wait \\
PREFETCHPERF & 统计 prefetch start、hit、miss、stall 等事件 \\
COLLECTPERF & 统计 collector packet、partial write、final write 和 context stall \\
PASSPERF & 统计 pass start 到 first fire、fire span、compute idle 等 \\
\bottomrule
\end{tabular}
\end{table}

这些计数器使本文能够说明每次优化后的瓶颈迁移。例如，drainpipe 主要降低 drain 阶段；continuous collector 降低 partial-PSUM 存储开销；backend full-tile 减少 spatial tile 固定成本；during-compute prefetch 降低下一 pass 准备造成的 compute idle。第六章将基于这些计数器进行实验分析。

\section{探索性 Column-Level PSUM}

当前 continuous PSUM collector 仍以完整 packet 为反馈粒度：collector 等待所有列 FIFO 都有数据后写入一个 $COLS\times2$ 通道 packet。Column-level PSUM 的探索方向是按列独立写入和读取 partial sum，使下一 pass 能够按照 systolic array 的列传播延迟更早获得可用数据。

本文已经实现了 column-level ping-pong buffer 和 column stream feeder 的基础模块，并完成模块级验证。但该路径尚未作为最终推荐板级配置，因此本文将其作为未来工作讨论，不将其计入核心性能贡献。这样可以避免把尚未完整端到端验证的实验性设计包装为成熟成果。

\section{本章小结}

本章从 partial-PSUM 反馈角度解释了后端大层的性能瓶颈，并介绍了 drainpipe、early drain、continuous collector、pass prefetch 和 during-compute prefetch 等调度优化。与第四章的数据搬运优化结合后，系统能够同时减少软件打包、DMA 启动、tile 边界、PSUM drain 和 pass 准备空泡，为第六章的端到端性能提升奠定基础。
